\section{Contexte et objectifs du projet}

\subsection{Contexte opérationnel}

L'exécution d'une dépense budgétaire repose sur une chaîne procédurale réglementée qui mobilise plusieurs intervenants internes et externes : l'ordonnateur, les fournisseurs, le contrôleur financier, le comptable matière et le comptable public. Cette organisation s'appuie sur le principe de séparation des fonctions d'initiative et d'engagement, de contrôle budgétaire préalable, de réception physique des biens et de prise en charge comptable.

En l'absence d'un outil fédérateur, la multiplicité des intervenants et des pièces justificatives expose le traitement des dossiers à des risques de lenteur dans la transmission de l'information, de perte de visibilité sur l'état d'avancement des demandes et de ressaisies multiples.

\subsection{Objectifs de l'application}

Le projet vise à concevoir une application informatique dédiée au suivi du processus de dépense budgétaire, depuis l'expression initiale du besoin jusqu'à la transmission du dossier d'ordonnancement au comptable public.

Trois objectifs opérationnels structurent cette démarche :

\begin{enumerate}
  \item \textbf{Structurer les interventions des différents acteurs :} attribuer à chaque intervenant un périmètre d'action clairement délimité dans le système, conforme à ses attributions administratives et financières.
  \item \textbf{Faciliter la circulation des dossiers :} fluidifier la transmission des projets d'actes, des informations de livraison, des factures, des pièces justificatives et des dossiers d'ordonnancement entre les acteurs successifs.
  \item \textbf{Assurer le suivi et la traçabilité des opérations :} permettre le suivi de l'évolution des demandes et conserver la trace des contrôles, visas, réceptions et signatures intervenus tout au long du cycle de vie de la dépense.
\end{enumerate}

\subsection{Limites de la présente phase d'étude}

Ce document s'inscrit dans une phase d'étude préliminaire. À ce stade :
\begin{itemize}
  \item il ne définit pas encore les objets métier détaillés, leurs attributs ou le modèle de données ;
  \item il ne préjuge d'aucun choix d'architecture logicielle ou d'infrastructure technique ;
  \item il s'arrête, conformément aux éléments recueillis, à la réception du dossier d'ordonnancement par le comptable public, sans couvrir les opérations ultérieures de règlement effectif ou de tenue de la comptabilité générale.
\end{itemize}
